% from ega4-2-fr.tex
\oldpage[IV-2]{99}

$\mathfrak p_i$ ($1\leq i\leq n$) sont les idéaux premiers minimaux de $A$,
il revient au même de dire que $A$ est universellement caténaire ou que
chacun des $A/\mathfrak p_i$ l'est. Cela montre aussi qu'il revient au même
de dire qu'un préschéma localement noethérien $X$ est universellement
caténaire, ou que $X_{\mathrm{red}}$ l'est.
\item[(iv)] Le critère (5.6.1, b)) et la remarque (i) montrent que, si $A$
est un anneau universellement caténaire, il en est de même de toute
$A$-algèbre essentiellement de type fini (1.3.8).
\end{enumerate}
\end{remarks}

\begin{proposition}[5.6.4]
\label{IV.2.s5.6.4-fr}
Tout anneau quotient d'un anneau régulier est universellement caténaire.
\end{proposition}

Tout revient à voir qu'un anneau régulier $A$ est universellement caténaire
(5.6.3, (iv)); or, on sait que $A[T_1,\ldots,T_n]$ est régulier (0, 17.3.7),
donc caténaire (0, 16.5.12), et l'on conclut par (5.6.1, a)).

En particulier, il résulte du théorème de Cohen (0, 19.8.8) que tout anneau
local noethérien complet est universellement caténaire. De même, toute
algèbre de type fini sur un corps étant quotient d'un anneau régulier, tout
préschéma localement de type fini sur un corps est universellement caténaire.

Nous verrons plus loin (6.3.7) que dans (5.6.4), on peut remplacer « anneau
régulier » par « anneau de Cohen-Macaulay ».

\begin{proposition}[5.6.5]
\EGAReaderTarget{ega4.c04.prop.5-6-5}{IV.2.s5.6.5-fr}
Soient $Y$ un préschéma irréductible localement noethérien, $X$ un préschéma
irréductible, $f:X\to Y$ un morphisme dominant localement de type fini. Soit
$\xi$ (resp. $\eta$) le point générique de $X$ (resp. $Y$), et soit
$e=\dim_\eta(f^{-1}(\eta))=\operatorname{deg.tr.}_{k(\eta)}k(\xi))$
(« dimension de la fibre générique », cf. $(0_{\mathrm I},2.1.8)$ et
(4.1.1)). Pour tout $x\in X$, on a alors, en posant $y=f(x)$,
\begin{equation}
e+\dim(\mathcal O_y)
\geq\operatorname{deg.tr.}_{k(y)}k(x)+\dim(\mathcal O_x)
\tag{5.6.5.1}\label{IV.2.s5.6.5.1-fr}
\end{equation}
\begin{equation}
\dim(\mathcal O_x)
\leq\dim(\mathcal O_y)
+\dim(\mathcal O_x\otimes_{\mathcal O_y}k(y))-\delta(x)
\tag{5.6.5.2}\label{IV.2.s5.6.5.2-fr}
\end{equation}
en posant $\delta(x)=\dim_x(f^{-1}(y))-e$.

En outre, si $Y$ est universellement caténaire, les deux membres de
(5.6.5.1) (resp. (5.6.5.2)) sont égaux. Si de plus $x$ est fermé dans
$f^{-1}(y)$, on a
\begin{equation}
\dim(\mathcal O_x)=\dim(\mathcal O_y)+e.
\tag{5.6.5.3}\label{IV.2.s5.6.5.3-fr}
\end{equation}
\end{proposition}

On peut évidemment se borner au cas où $X$ et $Y$ sont affines (donc $f$ de
type fini) et réduits (1.5.4), donc intègres. Comme $f$ est dominant,
$\mathcal O_y$ s'identifie alors à un sous-anneau de $\mathcal O_x$, et les
corps des fractions respectifs de $\mathcal O_y$ et de $\mathcal O_x$ à
$k(\eta)$ et $k(\xi)$ (I, 8.2.7); en outre, $\mathcal O_x$ est un anneau
local d'une $\mathcal O_y$-algèbre intègre de type fini $B$ en un idéal
premier $\mathfrak q$ (I, 3.6.5); en appliquant (5.5.8.1) à
$A=\mathcal O_y$, $B$ et $\mathfrak q$, on obtient (5.6.5.1). En outre,
$\mathcal O_x\otimes_{\mathcal O_y}k(y)=\mathcal O_x/\mathfrak m_y\mathcal O_x$
n'est autre que l'anneau local de la fibre $f^{-1}(y)$ au point $x$
(I, 3.6.1); comme $f^{-1}(y)$ est un préschéma de type fini sur $k(y)$, il
résulte de (5.2.3) que l'on a
\[
\dim(\mathcal O_x\otimes_{\mathcal O_y}k(y))
=\dim_x(f^{-1}(y))-\operatorname{deg.tr.}_{k(y)}k(x);
\]
l'inégalité (5.6.5.2) n'est donc qu'une autre forme de (5.6.5.1).

Le fait que les deux membres de (5.6.5.1) sont égaux lorsque $Y$ est
universellement caténaire n'est autre que l'égalité (5.6.1.1), appliquée à
$A=\mathcal O_y$ et $B'=\mathcal O_x$. Pour démontrer que l'on a de plus
(5.6.5.3) lorsque $x$ est fermé dans $f^{-1}(y)$, il suffit de remarquer que,
de façon générale, $\operatorname{deg.tr.}_{k(y)}k(x)$ est la dimension de
$\overline{\{x\}}\cap f^{-1}(y)$, comme il résulte de (5.2.1) appliqué au
sous-préschéma fermé réduit de $f^{-1}(y)$ ayant ce sous-espace comme espace
sous-jacent.

